% GENERATED FILE. Edit canonical lessons, then run npm run generate:books.
% canonical-source-hash: fnv1a64:bafd1405fe2712b6
% canonical-lessons: ES-C737-codo, ES-C737-muneca, ES-C737-cintura, ES-C737-cadera, ES-C737-costilla

\chapter{Del codo a la cadera}
\label{ch:es-del-codo-a-la-cadera}
\hlchaptermodality{737}

\begin{chapteropening}
\textbf{By the end of this chapter:} \emph{I can say where it hurts: elbow, wrist, waist, hip or ribs.}

Complete the last lesson of chapter 737: I can say where it hurts: elbow, wrist, waist, hip or ribs.
\end{chapteropening}

\section[el codo]{el codo — the elbow}

\label{lesson:ES-C737-codo}



\begin{quote}
Before the new one: say the Spanish for the beard, then the Spanish for the ear.
\end{quote}

\subsection*{You'll want to know: el codo}
\textbf{el codo} — \textquotedblleft{}the elbow\textquotedblright{}.

\begin{quote}
Me di un golpe en el codo.
\end{quote}

\begin{quote}
\emph{I banged my elbow.}
\end{quote}

\emph{Codo con codo} — side by side.

\begin{cousinweb}
\emph{codo} is inherited from Latin \emph{cubitus} \textquotedblleft{}elbow\textquotedblright{} (EtymDB), through Old Spanish \emph{cobdo}.
\end{cousinweb}

\begin{grammarlens}[title={What you've built}]
One more word for telling, explaining and describing at B1.
\end{grammarlens}

\subsection*{Your turn}
\begin{itemize}
\raggedright
  \item \emph{Say it:} \emph{el codo}
  \item \emph{Say it:} \emph{el codo}, once more
  \item \emph{Say it:} a sentence of your own with \emph{el codo}
  \item \emph{Recall:} say the Spanish for the beard, then the Spanish for the ear, then say \emph{el codo} again
\end{itemize}

\begin{tcolorbox}[breakable,colback=teal!4,colframe=teal!35!black,title={Before you move on}]
What does el codo mean? (\textquotedblleft{}The elbow\textquotedblright{}.) Say it once more.
\end{tcolorbox}
\section[la muñeca]{la muñeca — the wrist; the doll}

\label{lesson:ES-C737-muneca}



\begin{quote}
Before the new one: say the Spanish for the ear, then the Spanish for the elbow.
\end{quote}

\subsection*{You'll want to know: la muñeca}
\textbf{la muñeca} — \textquotedblleft{}the wrist; the doll\textquotedblright{}.

\begin{quote}
Me duele la muñeca de tanto escribir.
\end{quote}

\begin{quote}
\emph{My wrist hurts from so much writing.}
\end{quote}

Two senses: the wrist and a doll.

\begin{cousinweb}
For \emph{muñeca} the origin is not settled. Old Spanish \emph{munneca} meant a doll (EtymDB); the wrist sense may come from the knob of bone, but no source is sure.
\end{cousinweb}

\begin{grammarlens}[title={What you've built}]
One more word for telling, explaining and describing at B1.
\end{grammarlens}

\subsection*{Your turn}
\begin{itemize}
\raggedright
  \item \emph{Say it:} \emph{la muñeca}
  \item \emph{Say it:} \emph{la muñeca}, once more
  \item \emph{Say it:} a sentence of your own with \emph{la muñeca}
  \item \emph{Recall:} say the Spanish for the ear, then the Spanish for the elbow, then say \emph{la muñeca} again
\end{itemize}

\begin{tcolorbox}[breakable,colback=teal!4,colframe=teal!35!black,title={Before you move on}]
What does la muñeca mean? (\textquotedblleft{}The wrist; the doll\textquotedblright{}.) Say it once more.
\end{tcolorbox}
\section[la cintura]{la cintura — the waist}

\label{lesson:ES-C737-cintura}



\begin{quote}
Before the new one: say the Spanish for the elbow, then the Spanish for the wrist.
\end{quote}

\subsection*{You'll want to know: la cintura}
\textbf{la cintura} — \textquotedblleft{}the waist\textquotedblright{}.

\begin{quote}
El pantalón me queda grande de cintura.
\end{quote}

\begin{quote}
\emph{The trousers are too big at the waist.}
\end{quote}

\begin{cousinweb}
\emph{cintura} comes from Latin \emph{cinctūra} \textquotedblleft{}a girdle, a belt\textquotedblright{} (EtymDB), from \emph{cingere} \textquotedblleft{}to gird\textquotedblright{}.
\end{cousinweb}

\begin{grammarlens}[title={What you've built}]
One more word for telling, explaining and describing at B1.
\end{grammarlens}

\subsection*{Your turn}
\begin{itemize}
\raggedright
  \item \emph{Say it:} \emph{la cintura}
  \item \emph{Say it:} \emph{la cintura}, once more
  \item \emph{Say it:} a sentence of your own with \emph{la cintura}
  \item \emph{Recall:} say the Spanish for the elbow, then the Spanish for the wrist, then say \emph{la cintura} again
\end{itemize}

\begin{tcolorbox}[breakable,colback=teal!4,colframe=teal!35!black,title={Before you move on}]
What does la cintura mean? (\textquotedblleft{}The waist\textquotedblright{}.) Say it once more.
\end{tcolorbox}
\section[la cadera]{la cadera — the hip}

\label{lesson:ES-C737-cadera}



\begin{quote}
Before the new one: say the Spanish for the wrist, then the Spanish for the waist.
\end{quote}

\subsection*{You'll want to know: la cadera}
\textbf{la cadera} — \textquotedblleft{}the hip\textquotedblright{}.

\begin{quote}
Mi abuela se rompió la cadera.
\end{quote}

\begin{quote}
\emph{My grandmother broke her hip.}
\end{quote}

\begin{cousinweb}
\emph{cadera} is inherited from Latin \emph{cathedra} \textquotedblleft{}a seat, a chair\textquotedblright{}, through spoken \emph{cathegra}; the Latin came from Greek \emph{kathédra} (EtymDB). The hips are what you sit on.
\end{cousinweb}

\begin{grammarlens}[title={What you've built}]
One more word for telling, explaining and describing at B1.
\end{grammarlens}

\subsection*{Your turn}
\begin{itemize}
\raggedright
  \item \emph{Say it:} \emph{la cadera}
  \item \emph{Say it:} \emph{la cadera}, once more
  \item \emph{Say it:} a sentence of your own with \emph{la cadera}
  \item \emph{Recall:} say the Spanish for the wrist, then the Spanish for the waist, then say \emph{la cadera} again
\end{itemize}

\begin{tcolorbox}[breakable,colback=teal!4,colframe=teal!35!black,title={Before you move on}]
What does la cadera mean? (\textquotedblleft{}The hip\textquotedblright{}.) Say it once more.
\end{tcolorbox}
\section[la costilla]{la costilla — the rib}

\label{lesson:ES-C737-costilla}



\begin{quote}
Before the new one: say the Spanish for the waist, then the Spanish for the hip.
\end{quote}

\subsection*{You'll want to know: la costilla}
\textbf{la costilla} — \textquotedblleft{}the rib\textquotedblright{}.

\begin{quote}
Me duele una costilla al respirar.
\end{quote}

\begin{quote}
\emph{One of my ribs hurts when I breathe.}
\end{quote}

\emph{Las costillas de cerdo} — pork ribs.

\begin{cousinweb}
\emph{costilla} comes from Late Latin \emph{costula} \textquotedblleft{}little rib\textquotedblright{} (EtymDB), a small form of \emph{costa} \textquotedblleft{}rib, side\textquotedblright{}, the source of \emph{costa} \textquotedblleft{}coast\textquotedblright{}.
\end{cousinweb}

\begin{grammarlens}[title={What you've built}]
One more word for telling, explaining and describing at B1.
\end{grammarlens}

\subsection*{Your turn}
\begin{itemize}
\raggedright
  \item \emph{Say it:} \emph{la costilla}
  \item \emph{Say it:} \emph{la costilla}, once more
  \item \emph{Say it:} a sentence of your own with \emph{la costilla}
  \item \emph{Recall:} say the Spanish for the waist, then the Spanish for the hip, then say \emph{la costilla} again
\end{itemize}

\begin{tcolorbox}[breakable,colback=teal!4,colframe=teal!35!black,title={Before you move on}]
What does la costilla mean? (\textquotedblleft{}The rib\textquotedblright{}.) Say it once more.
\end{tcolorbox}
